\documentclass[11pt]{article}
\usepackage[margin=1in]{geometry}
\usepackage{amsmath, amssymb}
\usepackage{booktabs}
\usepackage{graphicx}
\usepackage{caption}
\usepackage{subcaption}
\usepackage{hyperref}
\usepackage{array}
\usepackage{multirow}
\usepackage{xcolor}
\usepackage{parskip}
\usepackage{enumitem}

\title{Ablation Study Plan: \\
       Two-Stage CKME Framework}
\author{}
\date{\today}

% Placeholder color for unfilled results
\newcommand{\tbd}{\textcolor{gray}{---}}

\begin{document}
\maketitle

\tableofcontents
\newpage

% ============================================================
\section{Overview}
% ============================================================

This document records the ablation study design for the two-stage CKME
(Conditional Kernel Mean Embedding) framework.
The full model (\textbf{CKME-full}) consists of:
\begin{enumerate}[noitemsep]
  \item \textbf{Stage 1}: Train CKMEModel on $n_0$ sites $\times$ $r_0$ reps,
        with cross-validated hyperparameters $(\ell_x, \lambda, h)$.
  \item \textbf{$S^0$ score}: Estimate uncertainty at candidate sites as
        $S^0(x) = \hat{q}_{1-\alpha/2}(x) - \hat{q}_{\alpha/2}(x)$
        (predicted interval width from Stage~1).
  \item \textbf{Stage 2}: Select $n_1$ new sites guided by $S^0$ scores
        (method: \texttt{mixed}); collect $r_1$ reps; calibrate split-CP.
  \item \textbf{Smooth indicator}: Logistic function with bandwidth $h > 0$
        for CDF estimation.
\end{enumerate}

Each ablation removes or replaces exactly one component while holding the rest
fixed.  All ablations use:
\begin{itemize}[noitemsep]
  \item The same macro-rep seeds.
  \item The same total sample budget
        $B = n_0 r_0 + n_1 r_1$.
  \item Evaluation metrics: marginal coverage, mean interval width,
        and interval score (IS).
  \item Simulators: six non-Gaussian cases
        (A1-S/L, B2-S/L, C1-S/L).
\end{itemize}

The interval score (Winkler score) is
\[
  \mathrm{IS}(L, U, Y) = (U - L)
  + \frac{2}{\alpha}(L - Y)_+
  + \frac{2}{\alpha}(Y - U)_+,
\]
where $(z)_+ = \max(z, 0)$.

% ============================================================
\section{Ablation 1: $S^0$-Guided Site Selection vs.\ Uniform LHS}
\label{sec:abl1}
% ============================================================

\subsection*{Research question}
Does the $S^0$ uncertainty score actually improve Stage~2 site selection,
or is space-filling Latin Hypercube Sampling (LHS) sufficient?

\subsection*{Variants}

\begin{center}
\begin{tabular}{lll}
  \toprule
  Label & Site selection & Config change \\
  \midrule
  \textbf{CKME-full} & $S^0$-guided (\texttt{method=mixed}) & baseline \\
  \textbf{CKME-no-S0} & Pure LHS (\texttt{method=lhs}) & \texttt{method=lhs} \\
  \bottomrule
\end{tabular}
\end{center}

\subsection*{Expected outcome}
CKME-full should yield narrower intervals at similar coverage,
especially in heteroscedastic regions where $S^0$ concentrates samples.

\subsection*{Status}
\textcolor{orange}{PLANNED} --- experiment not yet run.

\subsection*{Results}
\textit{(Fill in after experiment.)}

% ============================================================
\section{Ablation 2: Two-Stage vs.\ One-Stage}
\label{sec:abl2}
% ============================================================

\subsection*{Research question}
Is the two-stage structure necessary?
Does adaptive Stage~2 data collection improve CP calibration
over using Stage~1 data alone (with the same total budget)?

\subsection*{Variants}

\begin{center}
\begin{tabular}{lll}
  \toprule
  Label & Structure & Budget \\
  \midrule
  \textbf{CKME-full} & Stage 1 ($n_0, r_0$) + Stage 2 ($n_1, r_1$) & $B$ \\
  \textbf{CKME-1stage} & Stage 1 only, larger $n_0'$ or $r_0'$ & $B$ \\
  \bottomrule
\end{tabular}
\end{center}

Budget matching: $n_0' r_0' = n_0 r_0 + n_1 r_1 = B$.
CP is calibrated on the single Stage~1 dataset.

\subsection*{Expected outcome}
Two-stage should improve coverage in tail/heteroscedastic regions,
because Stage~2 concentrates calibration samples where coverage is hardest.
One-stage CP calibrated on uniform Stage~1 data may under-cover at the tails.

\subsection*{Status}
\textcolor{orange}{PLANNED} --- needs a one-stage runner script.

\subsection*{Results}
\textit{(Fill in after experiment.)}

% ============================================================
\section{Ablation 3: Smooth Indicator vs.\ Hard Step Function}
\label{sec:abl3}
% ============================================================

\subsection*{Research question}
Does the smooth logistic indicator (bandwidth $h > 0$) matter for
CDF estimation quality, compared to the hard Heaviside step function?

\subsection*{Variants}

\begin{center}
\begin{tabular}{lll}
  \toprule
  Label & Indicator & Bandwidth \\
  \midrule
  \textbf{CKME-smooth} & Logistic: $\sigma_h(t - y)$ & $h$ tuned by CV \\
  \textbf{CKME-hard} & Heaviside: $\mathbf{1}[t \ge y]$ & $h \to 0$ \\
  \bottomrule
\end{tabular}
\end{center}

\noindent
The smooth indicator is
\[
  \phi_h(t, y) = \frac{1}{1 + \exp\!\bigl(-(t - y)/h\bigr)}.
\]

\subsection*{Note}
\texttt{exp\_conditional\_coverage} already explored this partially
for 1D cases; those results can be reused or cited here.

\subsection*{Expected outcome}
Smooth indicator reduces gradient noise in the CRPS loss,
enabling better CV tuning.
Hard step may perform comparably for large $n$ but degrade for small $n$.

\subsection*{Status}
\textcolor{orange}{PLANNED} --- partial data from \texttt{exp\_conditional\_coverage}.

\subsection*{Results}
\textit{(Fill in after experiment.)}

% ============================================================
\section{Ablation 4: CV Hyperparameter Tuning vs.\ Fixed Params}
\label{sec:abl4}
% ============================================================

\subsection*{Research question}
How much does cross-validated selection of $(\ell_x, \lambda, h)$
contribute, compared to using fixed default values?

\subsection*{Variants}

\begin{center}
\begin{tabular}{lll}
  \toprule
  Label & Hyperparameters & Method \\
  \midrule
  \textbf{CKME-CV} & $(\ell_x, \lambda, h)$ by $k$-fold CRPS CV & \texttt{pretrain\_params.py} \\
  \textbf{CKME-fixed} & $\ell_x = 1.0,\; \lambda = 0.01,\; h = 0.1$ & hardcoded \\
  \bottomrule
\end{tabular}
\end{center}

\subsection*{Expected outcome}
CV tuning matters more for strongly non-Gaussian simulators
(B2-L, C1-L) where default $h$ may over-smooth or under-smooth the CDF.

\subsection*{Status}
\textcolor{orange}{PLANNED}.

\subsection*{Results}
\textit{(Fill in after experiment.)}

% ============================================================
\section{Ablation 5: Budget Allocation Between Stages}
\label{sec:abl5}
% ============================================================

\subsection*{Research question}
Given a fixed total budget $B$, what is the optimal split between
Stage~1 and Stage~2?

\subsection*{Variants}

\begin{center}
\begin{tabular}{lrrrrc}
  \toprule
  Label & $n_0$ & $r_0$ & $n_1$ & $r_1$ & Total $B$ \\
  \midrule
  \textbf{Heavy Stage 1} & 400 & 20 & 250 & 10 & 10{,}500 \\
  \textbf{Balanced (current)} & 250 & 20 & 500 & 10 & 10{,}000 \\
  \textbf{Heavy Stage 2} & 150 & 20 & 700 & 10 & 10{,}000 \\
  \bottomrule
\end{tabular}
\end{center}

\subsection*{Expected outcome}
An optimal balance exists:
too little Stage~1 leads to a poor $S^0$ score and bad site selection;
too little Stage~2 yields insufficient CP calibration data.

\subsection*{Status}
\textcolor{orange}{PLANNED} --- exploratory; run last if time permits.

\subsection*{Results}
\textit{(Fill in after experiment.)}

% ============================================================
\section{Summary Table}
\label{sec:summary}
% ============================================================

\begin{table}[h]
\centering
\caption{Ablation results summary (averaged over 6 simulators and all macro-reps).
         \tbd{} = not yet run.}
\begin{tabular}{llcccl}
  \toprule
  Ablation & Variant & Coverage & Width & IS & vs.\ CKME-full \\
  \midrule
  \multirow{2}{*}{1 (S$^0$ score)}
    & CKME-full   & \tbd & \tbd & \tbd & baseline \\
    & CKME-no-S0  & \tbd & \tbd & \tbd & \tbd \\
  \addlinespace
  \multirow{2}{*}{2 (two-stage)}
    & CKME-full   & \tbd & \tbd & \tbd & baseline \\
    & CKME-1stage & \tbd & \tbd & \tbd & \tbd \\
  \addlinespace
  \multirow{2}{*}{3 (indicator)}
    & CKME-smooth & \tbd & \tbd & \tbd & baseline \\
    & CKME-hard   & \tbd & \tbd & \tbd & \tbd \\
  \addlinespace
  \multirow{2}{*}{4 (CV tuning)}
    & CKME-CV     & \tbd & \tbd & \tbd & baseline \\
    & CKME-fixed  & \tbd & \tbd & \tbd & \tbd \\
  \bottomrule
\end{tabular}
\end{table}

% ============================================================
\section{Priority Order}
% ============================================================

\begin{enumerate}
  \item \textbf{Ablation~\ref{sec:abl1}} ($S^0$ score) ---
        most directly tied to the core claim; only requires
        changing \texttt{method=lhs} in the config.
  \item \textbf{Ablation~\ref{sec:abl2}} (two-stage necessity) ---
        second most important; needs a new one-stage runner.
  \item \textbf{Ablation~\ref{sec:abl3}} (smooth indicator) ---
        partial data already exists from \texttt{exp\_conditional\_coverage}.
  \item \textbf{Ablation~\ref{sec:abl4}} (CV tuning) ---
        sanity check; lower priority.
  \item \textbf{Ablation~\ref{sec:abl5}} (budget allocation) ---
        exploratory; run only if time permits.
\end{enumerate}

\end{document}
